\chapter*{Abstract}
\label{chap:abstract}
\addcontentsline{toc}{chapter}{\nameref{chap:abstract}}

Video surveillance systems generate massive volumes of footage daily, yet extracting useful information from these recordings remains a significant challenge. Conventional video analytics tools are restricted to predefined object categories, making them unable to handle the diverse and unpredictable queries that arise in real-world scenarios such as security investigations and retail monitoring. Even when relevant footage exists, studies have shown that only a fraction of it translates into actionable evidence, largely due to the difficulty of locating specific moments within extensive archives.

This project presents an open-vocabulary video search system that enables users to search through video footage using natural language descriptions or reference images. Unlike traditional systems, it does not rely on fixed category labels. Instead, it employs a region-level retrieval pipeline that first segments individual objects in each frame using the Segment Anything Model (SAM), then encodes each region into a shared vision-language embedding space using CLIP. These embeddings are stored with their corresponding timestamps in a vector database, allowing users to retrieve relevant moments through cosine similarity search. For identity-based queries, the system integrates InsightFace to detect, recognize, and cluster faces across footage.

Prior work on SAM-CLIP Search has demonstrated that region-level encoding outperforms frame-level encoding, achieving 92\% retrieval accuracy compared to 85\% with CLIP alone on standard benchmarks. This project builds on that validated approach and extends it to video by incorporating temporal indexing, face recognition, and an interactive search interface. The system is evaluated by comparing region-level and frame-level retrieval on both standard datasets and real surveillance footage, using metrics including Recall@5, Precision@5, retrieval accuracy, and query speed.